% Draft replacement for \section{Notifications} in
% text/chapters/6_implementation.tex (currently lines 804-824).
%
% Goals:
%   - Frame the design as a partial refactor from one shared model
%     into three independent ones, motivated by diverging fields.
%   - Correct two factual issues in the current prose:
%       (a) Only NotificationProfile is created from a signal handler.
%           The three Notification kinds are created directly from the
%           social and users views.
%       (b) read_at is one timestamp per User, not per User per kind.


\section{Notifications}
This section covers how Notifications are modelled, emitted, and
delivered. The Backend code is under
\path{src/backend/musikk/notifications}, the Frontend code under
\path{src/frontend/musikk/src/features/notifications}. The class
diagram is shown in Figure~\ref{fig:cd_notifications} and the
Frontend component hierarchy in Figure~\ref{fig:fe_notifications}.

There are three Notification kinds: a Reply to a Publication the User
authored, a new Follower, and a new Chat message. The first design
used one shared \texttt{Notification} base with three concrete
subclasses through Django's multi-table inheritance, and each row
held an \texttt{is\_read} boolean. This worked well at the start,
because every kind had the same two operations, list and mark read,
and one query covered all of them, but as the kinds gained more fields of
their own, the shared base became hard to work with. A Reply points to two
Publications, a Follower to two Users, and a Chat message to a Chat
together with the message and the receiver. The columns almost do not
overlap, so a shared row was mostly nulls, and every read needed an
extra join into one of the concrete tables.

The code is therefore in the middle of a change toward three separate
models. Each kind is now its own \texttt{BaseModel} subclass with its
own Serializer and custom Manager, and creation happens at the call
site that caused the Notification, in the social and users views.
Each Manager overrides \texttt{create} to emit a
\texttt{notifications.changed} event on the receiver's per-User
WebSocket group, which invalidates the Notifications query on every
open tab. Read state was also made simpler. The per-row
\texttt{is\_read} boolean was replaced by a single
\texttt{NotificationProfile.read\_at} timestamp per User, set up on
User creation through a \texttt{post\_save} signal handler. The
Serializer computes \texttt{is\_read} for each row by comparing
\texttt{date\_added} against that timestamp, so "mark all as read"
becomes one field update.

What still stays shared is the HTTP surface seen by the Frontend. One
\texttt{/notifications} endpoint returns a response with three
buckets, a \texttt{PATCH} on the same path updates \texttt{read\_at},
and \texttt{DELETE /notifications/<uuid>} checks the three tables in
turn to find the matching row. This shape still suits the Frontend,
which fetches all kinds together for the badge and the list view, but
the differences between kinds are no longer hidden behind a common
base.
